\documentclass[10pt,a4paper]{amsart}
\usepackage[margin=19mm]{geometry}
\usepackage{amsmath,amssymb,mathtools}
\usepackage[scr=rsfs,cal=euler]{mathalfa}
\usepackage{xfrac}
\usepackage[unicode,hidelinks,bookmarks=false]{hyperref}
\newcommand{\ie}{i.e.\ }
\newcommand{\cf}{cf.\ }
\newcommand{\resp}{respectively\ }
\newcommand{\Subsection}{\subsection}
\providecommand{\nobreakdash}{\nobreak\hbox{-}}
\setlength{\emergencystretch}{2em}
\multlinegap0pt
\usepackage{bm}
\usepackage{bbm}
\usepackage{dsfont}
\usepackage{xspace}
\usepackage{cancel}
\usepackage{mathtools}
\usepackage{amsthm}
\makeatletter
\@ifundefined{thm}{\newtheorem{thm}{Theorem}}{}
\@ifundefined{defi}{\newtheorem{defi}[thm]{Definition}}{}
\makeatother
\DeclareMathOperator{\s}{span}
\newcommand{\LF}{\mathfrak{F}}
\newcommand{\La}{\mathfrak{a}}
\newcommand{\Lg}{\mathfrak{g}}
\newcommand{\Lh}{\mathfrak{h}}
\newcommand{\Lz}{\mathfrak{z}}
\newcommand{\ov}{\overline}
\newcommand{\pkt}{\mathbin{\raise0pt\hbox{$\scriptstyle\bullet$}}}
\newcommand{\ra}{\rightarrow}  
\renewcommand{\phi}{\varphi}
\title[Affine cohomology classes for filiform Lie algebras]{Affine cohomology classes for filiform Lie algebras}
\author{Dietrich Burde}
\date{Source: arXiv:2601.09466v1 (CC BY 4.0)}
\begin{document}
\maketitle

\noindent\emph{Source and licence.} Affine cohomology classes for filiform Lie algebras. Version arXiv:2601.09466v1 (CC BY 4.0), distributed under the Creative Commons Attribution 4.0 International licence, as recorded in the arXiv licence metadata for this exact version and shown on the abstract page https://arxiv.org/abs/2601.09466v1. Full rights notice: licenses/arxiv-2601.09466-CC-BY-4.0.txt.\par\smallskip

\noindent\textit{Editorial scope.} Counted unit: the Thm whose source label is \texttt{ext} in the article (source0.tex lines 630--638, source label \texttt{ext}), delivered together with the article's own complete original proof of it (source0.tex lines 640--681, one \texttt{proof} environment of 42 lines). No mathematics is written, repaired or re-derived: statement and proof are the article's own text line for line. Required context retained uncounted: lsa1 (lines 200--210); affine (lines 215--222); lsa2 (lines 219--221). Every definition, display and label of those lines is retained unchanged. Omitted from this package and not redistributed: 1--199, 211--214, 222--629, 639--639, 682--1039. Nothing inside a retained excerpt is omitted or altered: every retained body is delivered exactly as the article's lines are written, so no omission receipt of any kind accompanies this package. Citations: the retained excerpts carry no citation command at all, so no bibliography entry of the article is transcribed and no reference is redistributed. Reference adaptation: none was needed and none was made; every reference inside the delivered unit resolves inside it, so no unresolved reference or citation is produced. Printed slips kept literal: no printed word, symbol, hypothesis or number of the counted statement or of its proof was altered. Layout and class: the article's own class is \texttt{amsart} with packages including amssymb; the wrapper is the lane's pinned \texttt{amsart} editorial wrapper at 10pt on A4, which restates the theorem-like environments the retained text needs and adds the article's own macro definitions that the retained text needs (\texttt{\textbackslash LF}, \texttt{\textbackslash La}, \texttt{\textbackslash Lg}, \texttt{\textbackslash Lh}, \texttt{\textbackslash Lz}, \texttt{\textbackslash ov}, \texttt{\textbackslash phi}, \texttt{\textbackslash pkt}, \texttt{\textbackslash ra}, \texttt{\textbackslash s}), with extra wrapper packages (bm, bbm, dsfont, xspace, cancel, mathtools). The delivered numbering is the unit's own. Assets: no retained span contains a figure environment, a graphics-inclusion command, a PDF-page inclusion, a table or a TikZ picture; no image file is copied into this package, the package declares no asset dependency and no image file is redistributed with it. Licence: version arXiv:2601.09466v1 of this article is distributed under the Creative Commons Attribution 4.0 International licence, as recorded in the arXiv licence metadata for this exact version and shown on the abstract page https://arxiv.org/abs/2601.09466v1. A full rights notice is delivered at licenses/arxiv-2601.09466-CC-BY-4.0.txt. Characters: every retained excerpt is pure ASCII, so the delivered engine, XeLaTeX with its default Latin Modern text font, reports no missing character.
\begin{defi}
Let $A$ be a vector space
over a field $K$ and let $A \times A \rightarrow A,
(x,y) \mapsto x\cdot y$ be a $K$--bilinear product which satisfies
\begin{equation}\label{lsa1}
x\cdot (y\cdot z)-(x\cdot y)\cdot z= y\cdot (x\cdot z)-(y\cdot x)\cdot z
\end{equation}
for all $x,y,z \in A$. Then $A$ together with the product
is called {\it left-symmetric algebra} or
LSA. The product is also called left-symmetric.
\end{defi}

\begin{defi}\label{affine}
An {\it affine structure} or {LSA--structure} on a Lie algebra
$\Lg$ is a $K$--bilinear product $\Lg \times \Lg \rightarrow \Lg$
which is left-symmetric and satisfies
\begin{equation}\label{lsa2}
[x,y]=x\cdot y -y\cdot x
\end{equation}
\end{defi}

\begin{thm} \label{ext}
Let $\Lg\in\LF_n(K)$ and suppose that $\Lg$
has an extension
\begin{equation*}
0 \rightarrow \La \xrightarrow{\iota} \Lh \xrightarrow{\pi}
\Lg \rightarrow 0
\end{equation*}
with $\iota(\La)=\Lz(\Lh)$. Then $\Lg$ admits an affine structure.
\end{thm}  

\begin{proof}
The first step of the proof consists in showing that
we may assume $\Lh\in \LF_{n+1}(K)$. For this we
refer the reader to the extended version of this article.
Let $(f_1,\dots ,f_{n+1})$ be an adapted basis for $\Lh$ and let
$e_i:=f_i$ mod $\Lz (\Lh)$ for $i=1,\dots, n$.
Then $(e_1,\dots ,e_n)$ is an adapted basis of $\Lg$. Let
$\Lh_3=\s \{f_3,\ldots ,f_{n+1}\}$ and $\Lg_2=\s \{e_2,\ldots ,e_{n}\}$.
There is a uniquely determined linear map $\phi : \Lg \ra \Lh_3$
satisfying $\phi (x)=[f_1,\ov{x}]_{\Lh}$ for all $x\in \Lg$ where
$\ov{x}\in \Lh$ is any element with $\pi (\ov{x})=x$. The restriction
of $\phi$ to $\Lg_2$ is bijective since it is evidently injective. Denote
its inverse by $\psi : \Lh_3 \ra \Lg_2$. Now set for all $x,y\in \Lg$    
\begin{equation}
\label{lsa}
x\pkt y:=\psi ([\ov{x},\phi (y)]_{\Lh})
\end{equation}
The formula is well defined since $[\ov{x},\phi (y)]_{\Lh}
=[\ov{x},[f_1,\ov{y}]]\in \Lh_3$.
We will show that it satisfies conditions \eqref{lsa1} and \eqref{lsa2} of
Definition $\ref{affine}$ and hence defines an affine structure on $\Lg$:

\begin{align*}
x\pkt y-y\pkt x & = \psi([\ov{x},[f_1,\ov{y}]]-[\ov{y},[f_1,\ov{x}]]) \\
 & = \psi([\ov{y},[f_1,\ov{x}]]-[f_1,[\ov{y},\ov{x}]] -[\ov{y},
[f_1,\ov{x}]]) \\
 & = \psi([f_1,[\ov{x},\ov{y}]])=\psi([f_1,\ov{[x,y]}_{\Lg}])
     =\psi(\phi([x,y]_{\Lg})) \\
 & = [x,y]_{\Lg}
\end{align*}
where the brackets are taken in $\Lh$ if not otherwise denoted.
Using the identity $[f_1,\ov{\psi(w)}]=w$ for all $w\in \Lh_3$
and again the Jacobi identity we obtain for all $x,y,z\in \Lg$:

\begin{align*}
x\pkt (y\pkt z)-y\pkt (x\pkt z) & = \psi ([\ov{x},[\ov{y},\phi(z)]]-
[\ov{y},[\ov{x},\phi(z)]]) \\
 & = [x,y]_{\Lg}\pkt z \\
 & = (x\pkt y) \pkt z-(y\pkt x)\pkt z
\end{align*}

\end{proof}    

\end{document}
